\subsection{A sufficient condition for behavioral recovery}
\label{app:behavior-recovery}

We apply margin-based stability reasoning \citep{tsuzuku2018lipschitz}
to cached decoding, adding a sensitivity assumption to Eq.~\eqref{eq:cache-distance}.

\paragraph{One-step guarantee.}
Fix the model, last prompt token $u_*$, positions and original mask $M_0$.
For vocabulary $\mathcal V$, let $F(\mathbf z)\in\mathbb R^{|\mathcal V|}$
map initial prefix caches to next-token logits. Assume finite sensitivity
$\kappa$ over the compared states in $\Omega$. The infinity norm is the
largest absolute change in any token score:
\begin{equation}
\|F(\mathbf z)-F(\mathbf z')\|_\infty
\le\kappa\|\mathbf z-\mathbf z'\|_2,\qquad
\mathbf z,\mathbf z'\in\Omega.
\label{eq:cache-lipschitz}
\end{equation}
For the donor's unique greedy choice $y_D$, define its positive margin
\begin{equation}
m_D=F_{y_D}(\mathbf z_D)-\max_{v\ne y_D}F_v(\mathbf z_D)>0.
\label{eq:donor-margin}
\end{equation}
\emph{Proposition.} If $2\kappa\varepsilon_{SC}<m_D$, then greedy
decoding from $\mathbf z_{SC}$ selects $y_D$.
\emph{Proof.} Each logit changes by at most
$\kappa\varepsilon_{SC}$. Thus, for every $v\ne y_D$,
\begin{equation}
F_{y_D}(\mathbf z_{SC})-F_v(\mathbf z_{SC})
\ge m_D-2\kappa\varepsilon_{SC}>0.
\label{eq:margin-preservation}
\end{equation}
The factor $2$ allows the winning logit to decrease and a competitor
to increase. Strict inequality excludes ties. \hfill$\square$

\paragraph{Complete generation.}
Let the donor produce $N$ tokens including its stopping token.
$F_t^D(\mathbf z)$ processes $u_*$ and fixed donor history $y_{D,<t}$
from initial cache $\mathbf z$, recomputing all new caches.
Assume Eq.~\eqref{eq:cache-lipschitz} for $F_t^D$ with bound $\kappa_t$
and donor margin $m_{D,t}>0$. The bound includes all intervening cache updates:
\begin{equation}
\bigl[\ \forall t\le N:\ 2\kappa_t\varepsilon_{SC}<m_{D,t}\ \bigr]
\quad\Longrightarrow\quad \hat y_{SC}=\hat y_D.
\label{eq:sequence-recovery}
\end{equation}
The first token agrees by the one-step result. If earlier tokens agree,
decoding follows the history defining $F_t^D$, so token $t$ agrees.
Induction under the common stopping rule gives an identical sequence,
which is correct whenever the donor's answer is correct.

\paragraph{Scope.}
Common valid $\kappa_t$ over $D$, $S$ and $SC$ give $SC$ no larger error
bounds (Eq.~\eqref{eq:cache-distance}). Losing cancellation can increase
actual errors, so accuracy gains are not guaranteed. Different sequences
may answer correctly, making this condition only sufficient.
$SC+\mathrm{cut}$ changes the mask and prefix and is excluded.
Sensitivity bounds and itemwise margins remain unverified.
